\begin{figure}[htbp]
\centering
\begin{tikzpicture}[font=\scriptsize]

% ---- sensores y panel ----
\node[caja, text width=2.2cm, align=center, draw=ambar, line width=1pt] (sens) at (0,0)
  {\textbf{Sensores}\\ inalámbricos\\ a batería};

\node[caja, text width=2.2cm, align=center, draw=ambar, line width=1pt] (panel) at (3.5,0)
  {\textbf{Panel} o\\ comunicador};

\draw[flecha] (sens) -- node[nota, above, font=\scriptsize, align=center] {radio\\propietaria} (panel);

% ---- central receptora ----
\node[caja, text width=2.5cm, align=center, draw=azuloscuro, line width=1.1pt,
      fill=azuloscuro!5] (arc) at (12.6,0)
  {\textbf{Central}\\ \textbf{receptora}\\ (ARC / CMS)};

% ---- punto de bifurcacion: el panel elige transporte ----
\coordinate (bif) at (5.5,0);
\draw[line width=0.7pt, color=gristexto] (panel.east) -- (bif);
\draw[line width=0.7pt, color=gristexto] (5.5,1.55) -- (5.5,-1.55);

\coordinate (conv) at (11.1,0);
\draw[line width=0.7pt, color=gristexto] (11.1,1.55) -- (11.1,-1.55);
\draw[flecha] (conv) -- (arc.west);

% ---- las tres generaciones, con etiqueta centrada en el tramo ----
\draw[flecha] (5.5,1.55) -- (11.1,1.55);
\node[nota, above=0.6mm of {(8.3,1.55)}, font=\scriptsize]
  {línea telefónica fija: \textbf{DTMF}\ \ \textcolor{gristexto}{\itshape (histórico)}};

\draw[flecha] (5.5,0.45) -- (11.1,0.45);
\node[nota, above=0.6mm of {(8.3,0.45)}, font=\scriptsize]
  {red celular: \textbf{DTMF} sobre 2G/3G\ \ \textcolor{gristexto}{\itshape (en retirada)}};

\draw[flechagruesa] (5.5,-1.15) -- (11.1,-1.15);
\node[nota, above=0.6mm of {(8.3,-1.15)}, font=\scriptsize\color{verdeacento}]
  {\textbf{internet: SIA DC-09}\ \ \itshape (actual, el adoptado)};

% ---- nota final ----
\node[cajaclave, text width=13.2cm, fill=verdeacento!6, font=\scriptsize, align=left]
      at (6.3,-3.1)
  {\textbf{En los tres casos el mensaje de aplicación es el mismo: Contact ID.} La «comunicación
   celular» es un cambio de transporte, no de protocolo: los comunicadores convierten la marcación de
   línea fija en marcación sobre red móvil conservando el mensaje intacto.\\[3pt]
   Por eso el mensaje tiene \textbf{16 dígitos y ni uno más}: debía caber en tonos telefónicos sobre una
   línea analógica. \textbf{Y esa restricción de hace cuatro décadas es la que hoy lo hace apropiado
   para LoRa}, un canal igualmente estrecho.};

\end{tikzpicture}
\caption{Cadena de comunicación de una alarma monitoreada. El formato del mensaje se conserva a través
de tres generaciones de transporte.}
\label{fig:alarmas}
\end{figure}
